% Ad Familiares 10.25 --- headnote
% ad-familiares-10-25, 26 May 43 BC, Rome

\textit{Cicero to C. Furnius, from Rome on 26 May 43 BC ---
Perseus dateline \textit{Scr. Romae circ. vii K. Iun. a. 711
(43)}. Furnius is Plancus's legatus and political fixer in
the city, and the addressee here is the same Furnius whose
errand to Octavian is described in Plancus's letter of
late July (10.24, $\S$~4--7). At the time of writing he is
travelling between Plancus's camp in Gaul and Rome,
carrying letters and intentions in both directions, and
already eyeing the practor's elections.}

\textit{The whole letter is a single piece of advice, given
with the controlled emphasis Cicero reserves for friends
he respects. Furnius has been canvassing in absentia (or
preparing to come down for the elections in person); Cicero
tells him to stay where he is, to finish ``what is left of
the war'' alongside Plancus, and to set his standing
($\textit{dignitas}$) ahead of his ambition. The career
arithmetic in section 2 is precise: Furnius is not yet at
the statutory year for the praetorship --- he has not held
the aedileship, so no fixed two-year interval ties his
hand --- and the candidacy will be more brilliant under
Plancus's consulship in any case. The closing scene is
domestic and characteristic: Cicero has rehearsed this
counsel at home in the presence of Quintus, Caecina,
Calvisius, and Furnius's own freedman Dardanus, and
reports their assent --- ``but you yourself will judge
best.''}
